\documentclass{ctexart}
\usepackage{amsmath, amssymb}
\usepackage{geometry}
\geometry{a4paper, margin=2cm}

\title{\textbf{第10章 恒定磁场} --- 例题}
\author{}
\date{}

\begin{document}
\maketitle

\begin{enumerate}

\item（例题1）计算长为 $L$ 的载流直导线在均匀磁场 $B$ 中所受的力。

\item（例题2）计算载流曲线在均匀磁场 $B$ 中所受的力。

\item 在均匀磁场中放置一任意形状的导线，电流强度为 $I$，求此段载流导线受的磁力。

\item 求两平行无限长直导线之间的相互作用力。

\item（例9.13，p.390）无限长载流直导线 $I_1$ 沿半径为 $R$ 的圆形载流导线 $I_2$ 的直径 $AB$ 放置。求：(1) 半圆弧 $ACB$ 所受安培力的大小和方向；(2) 整个圆形电流所受安培力的大小和方向。

\item 一半径为 $R$、载有电流 $I$ 的线圈，放在均匀的外磁场 $B$ 中，求此线圈中的张力。

\item 求一载流导线框在无限长直导线磁场中的受力和运动趋势。

\item（例9.15，p.399）在匀强磁场中，有一半径为 $R$ 的半圆形平面载流线圈，通有电流 $I$，$B$ 的方向与线圈平面平行。求：(1) 线圈所受安培力对 $y$ 轴之力矩；(2) 线圈平面转过 $\pi/2$ 时，磁力矩 $M$ 所做的功。

\item 求绕轴旋转的带电圆盘轴线上的磁场和圆盘的磁矩（书中题9.6，p.373）。

\item（书中例题9.7，p.375）长 $L = 0.1\,\text{m}$，带电量 $q = 1\times10^{-10}\,\text{C}$ 的均匀带电细棒，以速度 $v = 1\,\text{m/s}$ 沿 $x$ 轴正方向运动。当细棒运动到与 $y$ 轴重合时，细棒下端与坐标原点 $O$ 的距离 $a = 0.1\,\text{m}$。求此时坐标原点处磁感应强度 $B$ 的大小。

\item 如图的导线，已知电荷线密度为 $\lambda$，当绕 $O$ 点以 $\omega$ 转动时，求 $O$ 点的磁感应强度。

\item（例题1）在磁感应强度为 $B$ 的均匀磁场中，有一半径为 $R$ 的半球面 $S$，$S$ 的边线所在平面法线方向 $n$ 与 $B$ 的夹角为 $\alpha$，求通过半球面 $S$ 的磁通量。

\item（例题：9.18，p.419）无限长圆柱形铜线，外面包一层相对磁导率为 $\mu_r$ 的圆筒形磁介质。导线的半径为 $R_1$，磁介质的外半径为 $R_2$，铜线内均匀分布的电流 $I$，铜线的相对磁导率为 $1$。求无限长圆柱形铜线和介质内外的磁场强度 $H$ 和磁感应强度 $B$。

\item（例题2）无限长直导线通以电流 $I$，求通过如图所示的矩形面积的磁通量。

\item 一无限长载流直导线，其外包围一层磁介质，相对磁导率 $\mu_r > 1$。求：(1) 磁介质中的磁感应强度；(2) 介质内外界面上的束缚电流密度。

\item 求无限长圆柱面电流的磁场分布。

\item（例题9.9，p.384）无限长载流螺线管通有电流 $I$，单位长度上的匝数为 $n$。求螺线管内外的磁感应强度 $B$。

\item 求螺绕环电流的磁场分布及螺绕环内的磁通量。

\item 求无限大平面电流的磁场。

\end{enumerate}

\end{document}
